\chapter*{附录 C\quad 研读札记}
\addcontentsline{toc}{chapter}{附录 C\quad 研读札记}
\markboth{附录 C\quad 研读札记}{}

本附录是译校者本人的读书笔记，按章排列，与各章末的“本章评注”分工不同：本章评注面向读者概览全章，研读札记记录译校者自己在读这一章时的理解路径、抓到的技术要点、以及尚未解决的疑问。凡与译文有出入之处，一律以页图为准。

\section*{研读札记 00 \textperiodcentered{} 绪论（原书 p.1--8）}
\addcontentsline{toc}{section}{研读札记 00 \textperiodcentered{} 绪论（原书 p.1--8）}
\markboth{研读札记 00 \textperiodcentered{} 绪论（原书 p.1--8）}{}

\begin{quote}\small 这一组札记是我的读书笔记，不是译注，也不重复章末“本章评注”。它记录的是：这一章讲了什么、作者为什么这样安排、我在其中学到了什么、以及我还没弄明白的地方。凡与译文有出入之处，以页图为准。\end{quote}

\subsection*{一、这本绪论其实是一张"纲领地图"}

作者开篇不谈定理，而先讲一个问题：\textbf{能不能在每一个双有理等价类里挑出一个"最简单"的代表，然后只研究这些代表？} 这是分类学的思路，也是全书的问题意识。

\begin{itemize}
\item 曲线：每个不可约曲线都双有理等价于\textbf{唯一}一条非奇异射影曲线。所以曲线的情形是"完全解决"的典范——研究一条曲线 = 研究它所在的双有理等价类。
\item 曲面：一个不可约曲面双有理等价于\textbf{无穷多}条非奇异射影曲面。于是需要"极小模型理论"来挑代表：只要曲面上有自交数 $-1$ 的非奇异有理曲线，就可以把它收缩成一点，得到另一条非奇异射影曲面；反复收缩，"多数情形"下得到唯一的极小模型；少数情形下极小模型不唯一，但这些情形可以显式写清楚（第 1 章 \S1.4 会兑现这句话）。
\item 高维：这个办法一开始被认为要失败，直到 Reid 与森重文在 1980 年代初指出：\textbf{只要允许"某种温和的奇点"，高维极小模型理论是可能的}。这就是极小模型纲领（MMP / 森纲领）。1988 年森重文完成三维情形。
\end{itemize}

我在这里学到的第一件事是一个历史教训：\textbf{高维的困难不是"能不能收缩"，而是"收缩之后得到的簇不再是簇"}。允许奇点不是妥协，而是让理论闭合的必要条件。这也解释了全书为何花整整两章（第 4、5 章）专讲奇点分类——那不是预备知识，而是纲领的一部分。

\subsection*{二、全书结构的一句话版本}

作者自己给的路线图，我抄下来当索引：

\par\smallskip\noindent
\begin{longtable}{@{}>{\raggedright\arraybackslash}p{0.200\textwidth}>{\raggedright\arraybackslash}p{0.340\textwidth}>{\raggedright\arraybackslash}p{0.380\textwidth}@{}}
章 & 内容 & 在纲领中的位置 \\
\midrule
\endhead
1 & 有理曲线与标准类 & 用几何方法（Fano 簇、极射线）把纲领"看见"一遍 \\
2 & MMP 引论、奇点、消没定理 & 准备语言：klt/lc、极小模型、相对化 \\
3 & 锥定理 & 纲领能运行的\textbf{技术核心} \\
4 & 二维奇点（Corti 执笔） & 三维理论的局部模型 \\
5 & MMP 的奇点：终端、标准、伴随反转、双有理理论 & 纲领每一步的"路况" \\
6 & 三维翻模 & 收缩之后仍需改向时的新操作（翻模比翻转好懂） \\
7 & 半稳定翻转 & 三维翻转的存在性，经半稳定情形大幅简化 \\
\end{longtable}
\par\smallskip

两个"旁支"也值得注意：\S4.5（超曲面奇点的形变）与 \S5.5（双有理理论）。作者明说，这两节之所以写进来，是因为现有文献对它们覆盖不足，而书里恰好要用。\textbf{这句话本身就是一条方法论：教科书的价值有时不在主干，而在把散落在论文里的必需引理整理成可读的形式。}

\subsection*{三、我读到的三个"预埋点"}

\begin{enumerate}
\item \textbf{"-1 曲线收缩"是 Castelnuovo 定理的一半，反过来用}：曲面理论里它是判据（有 $-1$ 曲线 $\Longleftrightarrow$ 存在非平凡双有理态射）；在 MMP 里它是\textbf{算法步}（收缩它，继续走）。同一件事，从"判定"变成"操作"，整个几何就活了。
\item \textbf{有理曲线是"负性"的探测器}：第 1 章的核心洞察是，$K_X$ 的负性会强迫有理曲线出现，而有理曲线的缺失意味着双有理几何的刚性。这条线一直贯穿到锥定理：\textbf{极射线总被有理曲线生成}。
\item \textbf{记号体系即世界观}：绪论末尾的记号 0.4 有 17 条（(0)--(16)），其中真正影响全书的是：$\mathbb{Q}$-除子/$\mathbb{Q}$-Cartier/$\mathbb{Q}$-分解、除子的线性等价与数值等价、nef、$f$-等价（相对情形）、双有理变换 $f_*Z$、取整 $\lfloor D\rfloor,\lceil D\rceil,\{D\}$。\textbf{相对版本（$f$-nef、$f$-丰富、$f$-等价）不是装饰}——MMP 的每一步都是相对某个收缩态射做的，绝对版本只是 $f$ 为常值态射的特例。
\end{enumerate}

\subsection*{四、一个容易错过的细节：尾符三档}

绪论第 (16) 条定了全书的结束符号规则，而它比中文读者的习惯复杂一层：

\begin{itemize}
\item 证明末尾：$\blacksquare$；定义、例等末尾：$\square$；
\item 定理、命题等\textbf{不给出证明}时：$\blacksquare$；
\item \textbf{证明在后面给出}时：$\square$；
\item 还有第三档（原书没有单列，但正文里实际出现）：\textbf{证明紧接其后}时，声明末尾\textbf{什么都不印}，$\blacksquare$ 只在证明末尾出现。
\end{itemize}

这条规则对读证明型教材很重要：看到 $\square$ 却找不到证明，说明证明在后面某处，不能当成"未证"。我在翻译时把它写进了 preamble 的三档环境，并在正文里逐条核对页图——这条规则错了，全书的"可证性"地图就错了。

\subsection*{五、我的疑问与待办}

\begin{itemize}
\item 绪论提到"$(-1)$ 曲线收缩后得到另一条非奇异射影曲面"，但没提\textbf{收缩的唯一性}（Moishezon 的例子说明个别情形下不唯一）；第 1 章 \S1.4 应该补上，届时核对。
\item 日译本比英文原著多了 (0) 条（variety 默认既约、scheme 不假设），并补了日本教材（[永田 74]、[松村 80]、[宮西 90]）的指引。这是译者加的，不是 Kollár--Mori 原文——\textbf{凡属此类增补，我在译注里都要标明"日译本增补"}，避免读者误以为原著如此。
\item 绪论说"三维情形由森重文 1988 年完成"，但翻转存在性的证明细节要到第 7 章；我打算在第 7 章读完后回头检验这句话的分量。
\end{itemize}

\subsection*{六、如果只让我记一件事}

\textbf{双有理几何的分类问题，本质上是一个"如何在等价类中选代表"的问题；MMP 就是那套选代表的算法，而它必然把光滑流形带到带奇点的世界。} 全书后面所有技术，都是为了让这个算法在每一步都说得清楚。

\section*{研读札记 01 \textperiodcentered{} 第 1 章 有理曲线与曲线锥（原书 p.9--46）}
\addcontentsline{toc}{section}{研读札记 01 \textperiodcentered{} 第 1 章 有理曲线与曲线锥（原书 p.9--46）}
\markboth{研读札记 01 \textperiodcentered{} 第 1 章 有理曲线与曲线锥（原书 p.9--46）}{}

第 1 章是全书的"技术母机"：后面所有章都在用它在两种情形下造有理曲线的手法，以及它建立的曲线锥语言。我读这一章花的时间比读第 3 章还多，原因不是它难，而是它\textbf{同时给了两条互不相同的证明路线}，而我一开始把它们当成了同一条的两种写法。下面是按节顺序的笔记。

\subsection*{一、先说结论：这一章到底给了什么}

作者在 \S1.1 与 \S1.2 各证一次"负性迫使有理曲线存在"，两次的假设与工具完全不同：

\par\smallskip\noindent
\begin{longtable}{@{}>{\raggedright\arraybackslash}p{0.225\textwidth}>{\raggedright\arraybackslash}p{0.225\textwidth}>{\raggedright\arraybackslash}p{0.225\textwidth}>{\raggedright\arraybackslash}p{0.225\textwidth}@{}}
节 & 假设 & 工具 & 结论 \\
\midrule
\endhead
\S1.1 & $-K_X$ 丰富（Fano） & 双有理映射的不定点、形变 & 定理 1.2、定理 1.10：过任意点有有理曲线，且 $-(C\cdot K_X)\le\dim X+1$ \\
\S1.2 & 存在 $C'$ 使 $-(C'\cdot K_X)>0$ & 折弯引理（推论 1.7、引理 1.9）+ 正特征 Frobenius & 定理 1.13：有理曲线，且同时控制 $K_X$-次数与 $H$-次数之比 \\
\end{longtable}
\par\smallskip

我原来以为 \S1.2 是 \S1.1 的推广（"把 $-K_X$ 丰富放宽为某条曲线为负"），读下来才明白：\textbf{\S1.2 的技术是新的，\S1.1 反而更初等}。\S1.1 靠的是"双有理映射若有不定点，像中必含有理曲线"（命题 1.3、推论 1.4/1.5）与例 1.8 划出的边界：紧复流形上这些结论不成立，代数性不可省。\S1.2 才是把"负性"真正变成一个可操作的形变论证。

\subsection*{二、折弯引理：我抓到的那个"小机关"}

折弯引理（推论 1.7、引理 1.9）的证明骨架是：取一条曲线 $C$，\textbf{固定它的一点}去看形变空间，形变空间的维数估计是

\[
\dim \ge -(C\cdot f^*K_X)+(1-g(C))\dim X .
\]

只要右边 $\ge 0$，$C$ 就能在这样的形变中退化，退化后分解成若干曲线之和，其中至少一条的 $K_X$-次数严格更小（"折弯"）。反复折弯，次数降到不能降，就得到有理曲线。

我原来把这一步想成"曲线在曲面里滑动"，其实关键在于\textbf{两个项的符号}：

\begin{itemize}
\item $-(C\cdot K_X)>0$ 是假设给的，提供正的维数；
\item $(1-g(C))\dim X$ 在 $g(C)\ge1$ 时\textbf{非正}，会吃掉维数。
\end{itemize}

所以亏格一高就完全失效——这正是为什么结论必须是"有理曲线"而不能是"低亏格曲线"。也正因如此，$\dim X$ 因子越大越有利，这解释了后面所有高维论断里 $\dim X$ 出现在不等式右边的写法。定理 1.10 的界 $-(C\cdot K_X)\le\dim X+1$ 也是同一套估计的直接产物。

\subsection*{三、特征 $p$ 化：本章最值得学的技术}

\S1.2 最后要的是"高次自同态"造出的曲线族，而特征 $0$ 下没有现成的自同态。Kollár 的做法是把整个构造搬到特征 $p$：用 Frobenius 态射的 $m$ 次幂造出 $p^m$-次自同态，再用 generic flatness 保证交点数 $-(C\cdot K_X)$ 与亏格 $g(C)$ 在"几乎所有素数"上不随 $p$ 变化，于是这些素数上都造得出有理曲线；最后由原理 1.11 下降回特征 $0$。

我在这里停留很久，因为这是一种\textbf{我此前没用过的推理形态}：不是"在特征 $0$ 里做近似"，而是"在无穷多个特征 $p$ 的模型上做精确构造，再用有限生成性把它们收回到一般纤维"。技术要点有三条，我抄在下面：

\begin{enumerate}
\item 一切对象先在一族有限生成 $\Z$-概形上实现（"spreading out"），这样才有"对几乎所有素数"可言；
\item generic flatness 保证纤维上的 Hilbert 多项式（从而交点数、亏格）稳定；
\item 原理 1.11 是"下降"那一步：在几乎所有纤维上成立的构造，回到特征 $0$ 的一般纤维仍成立。
\end{enumerate}

\textbf{这条路线与第 3 章的无基点方法是两回事}。第 3 章用消没定理 + 无基点定理（非消没定理）造曲线，一步都不提特征 $p$。我原来以为第 3 章只是"把第 1 章做得更一般"，其实第 3 章是\textbf{为了做到 klt 对（有奇点）的情形，必须换掉整个工具链}，因为折弯引理依赖非奇异与形变空间的维数估计，在有奇点的对上根本无从谈起。看清这一点，第 3 章开头那串"非消没$\Rightarrow$无基点$\Rightarrow$有理性$\Rightarrow$锥定理"的推导链才显得合理。

\subsection*{四、曲线锥与收缩：从"造曲线"到"造算法"}

\S1.3 把前面造出的曲线收进一个几何对象：定义 1.15--1.17 给出数值等价、$NE(X)$ 与其闭包 $\overline{NE}(X)$。定理 1.24（锥定理）说 $\overline{NE}(X)$ 由 $K_X\ge0$ 的部分与可数条极射线生成；把边界微移到 $(K_X+\epsilon H)\ge0$ 后只剩有限条——这一步（推论 1.19 的紧性）是"算法化"的关键，因为极小模型纲领每一步都要在\textbf{有限}多条极射线里挑一条收缩。

我原来以为"极射线"是个抽象概念，读到定义 1.25（收缩）才明白它的用法：极射线对应一个收缩射 $f:X\to Y$，而 MMP 的一步就是"找 $K_X$-负的极射线，收缩它"。例 1.27 提醒收缩之后未必射影，这不是技术瑕疵而是真实现象。

\S1.4 在曲面上把这一步走到底：定理 1.28 列出三种收缩（到光滑曲线、到点、以及例外型），每收缩一次 Picard 数 $\rho$ 下降 $1$，因此必然终止，得到极小模型（定理 1.29、概要 1.31 的步骤 0--5）。三维（定理 1.32）出现新现象：出现小收缩，极小模型不一定存在，需要翻转——这是第 6 章的主题。\textbf{"$\rho$ 每步降 1"这个账本}是我在这一章学到的第二个可迁移的工具：终止性不来自"每步变小"，而来自一个整数值不变量严格下降。

\subsection*{五、丰富性判别法：两种正性之间的翻译}

\S1.5 我读得最吃力，因为它把前面所有东西反过来用。两个判别法其实是同一件事的两面：

\begin{itemize}
\item \textbf{Kleiman 判别法（定理 1.18）}：$D$ 丰富 $\iff$ $D$ 在 $\overline{NE}(X)\setminus\{0\}$ 上 $>0$。这是把丰富性翻译成\textbf{锥上的数值正性}；
\item \textbf{中井--Moishezon 判别法（定理 1.37）}：$L$ 丰富 $\iff$ 对一切闭子概形 $Z$ 有 $(L^{\dim Z}\cdot Z)>0$。这是把丰富性翻译成\textbf{逐个子概形的几何正性}。
\end{itemize}

两者之间的桥是定理 1.38（nef $\Rightarrow$ 交点非负，Kleiman 的原始结果）：它保证"锥上非负"可以一路推到"每个子概形上非负"，从而把 1.39 的锥判据接到 1.37 的子概形判据上。我在 ch01d 的译注里专门记了定理 1.38 证明里那个动作：$f(1,t)$ 的系数非负 $\Rightarrow$ 严格增；若 $(L^n)<0$ 就有唯一的 $t_0>0$ 使 $f(1,t_0)=0$，再把 $s>t_0$ 处的丰富性经不等式 (1.6) 与极限传回 $t_0$，得到矛盾。\textbf{"负的量无处可逃"}——这种用"一点点正性 + 极限"堵死的写法，是本章我第二个想记住的证明套路。

相对化（命题 1.41、定理 1.42、1.44、命题 1.45）用的是另一套：沿纤维方向对 $i$ 作降链归纳证明相对消没定理，再用中山引理把"逐点可生成"提升为"邻域上可生成"。例 1.46 是我在这一章最大的收获之一：$D=\Gamma_1-\Gamma_0$ 在 $E\times E$ 上每个纤维次数为 $0$、数值 $f$-平凡，却不是 nef，而且对任何 $A$ 都有 $(D+f^*A)^2=-2$——\textbf{f-nef 与 nef 之间没有"加个 $f^*A$ 就救回来"的余地}，所以相对判别法必须用 $\overline{NE}(X/Y)$ 而不能用 $\overline{NE}(X)$。

\subsection*{六、与我原有理解的出入}

\begin{enumerate}
\item 我原以为"锥定理"的证明必然用消没定理。第 1 章的锥定理（定理 1.24）用的是折弯引理与特征 $p$ 自同态，\textbf{一次消没都没用}；用消没的是第 3 章的版本。
\item 我原以为 nef 就是"对一切曲线非负"的数值条件。例 1.46 说明在相对情形下这个直觉会骗人：$f$-nef 只约束纤维方向，纤维外可以负到无法补救。
\item 我原以为丰富性判别法是"判据"，属于技术细节。实际上它是\textbf{把数值信息翻译成几何信息}的桥，第 3 章锥定理的"丰富 + 典范"配方（某 $aD-K_X-\Delta$ 丰富且大）正是这座桥的另一端。
\item 折弯引理我原以为靠"曲线能在曲面里动起来"。真正的条件是形变空间维数估计里 $(1-g(C))\dim X$ 这一项——亏格必须低，所以结论只能是有理曲线。
\end{enumerate}

\subsection*{七、未解疑问}

\begin{enumerate}
\item 定理 1.37 的证明只对射影的 $X$ 做，注 1.40 说固有非奇异（或 $\Q$-分解）时也成立而一般固有簇不成立。\textbf{"固有但非射影"这一档到底怎么构造反例？} 书中未给，我想在第 3 章之后回头补读 [Kle66]。
\item 特征 $p$ 化那一步我只跟到"对几乎所有素数构造成功"，原理 1.11 的下降机制（有限生成 $\Z$-概形 + 一般纤维）我自己还没有完整复述过一遍；这是我想动手验算的地方。
\item 例 1.23(2) 把 $\overline{NE}(X)$ 的截口说成"圆 $C$ 的圆周及其内部"，我在页图（图 1.4）上核对过，但"截口为什么正好是圆盘"这一点的几何来源书中没展开，我想在第 2 章例 2.16 的框架下再读一次。
\item 定理 1.32 只说三维出现小收缩与例外型收缩等新现象，\textbf{并未给出三维极小模型存在性的完整陈述}（那是第 6 章）。本章的"概要 1.31"只到曲面；三维的完整算法我在第 6 章回看。
\end{enumerate}

\section*{研读札记 02 \textperiodcentered{} 第 2 章 极小模型纲领引论（原书 p.47--94）}
\addcontentsline{toc}{section}{研读札记 02 \textperiodcentered{} 第 2 章 极小模型纲领引论（原书 p.47--94）}
\markboth{研读札记 02 \textperiodcentered{} 第 2 章 极小模型纲领引论（原书 p.47--94）}{}

本章是全书从“几何操作”转向“纲领语言”的转折点。第 1 章还在曲面上逐个收缩 $(-1)$-曲线，本章开始把整套操作抽象成“锥定理 + 收缩 + 翻转 + 终止”的四件套，并明确宣告：\textbf{三维以上必须在允许奇点的簇范畴里工作}。我的收获主要不在具体定理，而在看清了这个范畴为什么是这样被设计出来的。

\subsection*{一、本章在做什么}

\begin{itemize}
\item \S2.1 用例 2.1（三维商奇点 $\frac{1}{r}(1,a,r-a)$）说明非奇异簇的范畴对 MMP 不封闭：收缩一条极射线后得到的簇一般带奇点。于是纲领改在“终端奇点”这一类上运行。四步流程：$K_X$ nef 则停；否则取极射线并收缩；收缩若为除子型（divisorial）就继续，若为小收缩则作翻转；如此循环。
\item \S2.2 把纲领相对化，并推广到解析情形、带群作用、对数对 $(X,\Delta)$。四者共用同一套锥定理与收缩定理，因此框架可以整体照搬——这一点我认为是本书写作上最省力也最诚实的地方：先做非奇异情形，再声明“所有论证都相对化”，把重复劳动一次性卖掉。
\item \S2.3 是真正的技术核心：用差异系数 $a(E,X,\Delta)$ 统一地定义终端、标准、klt、dlt 四档，并证明 dlt 在消解、收缩、翻转下保持。
\item \S2.4--2.5 提供纲领的燃料：小平消没定理与 Kawamata--Viehweg 消没定理。没有消没，翻转的存在性与相对收缩都无从谈起。
\end{itemize}

\subsection*{二、我抓到的技术要点}

\begin{enumerate}
\item \textbf{差异系数是不变量，不是坐标量。} 定义 $K_Y=f^*(K_X+\Delta)+\sum a(E,X,\Delta)E$ 依赖于模型 $Y$，但 $a(E,X,\Delta)$ 只依赖 $E$ 对应的离散赋值 $v_E$（通过局部环 $\cO_{e,X}$ 与它的赋值环）。这正是 \S2.3 全部分类“可以只谈赋值、不谈模型”的根据。第 3 章证明“翻转前后奇点类别不变”时，用的就是同一个不变性。
\item \textbf{klt 与 dlt 的差别只在系数 $\ge1$ 与 $<1$。} klt 要求所有 $a(E,X,\Delta)>-1$；dlt 放宽为“存在 snc 的对数消解，且系数 $\le1$ 的边界成分正规”。放宽的代价是边界必须既约——这解释了为什么第 5 章要用整整一节去证明 dlt 对的有理性（定理 5.22），而 klt 情形几乎是免费的。
\item \textbf{大（big）与 nef 的组合是可加的。} 命题 2.61 把“nef 且大”等价于“丰富 $\Q$-除子 + 小有效除子之和”，这是 Kawamata--Viehweg 消没定理能推广到“接近丰富”的全部秘密：附加一个足够小的有效除子，不改变消没的结论，却把条件降到可验证的数值层面。
\item \textbf{消没定理的系数条件 $0\le a_i<1$ 不能碰 1。} 一旦某个 $a_i=1$，klt 性失效，反例立刻出现（圆锥上的 $\cO(-D)$）。这是贯穿全书的一条红线，第 3 章处理 lc 对时必须先扰动 $\Delta\mapsto(1-\delta)\Delta+\epsilon H$ 化为 klt，原因就在这里。
\item \textbf{相对化不是装饰。} 定理 2.64 的相对形式（推论 2.68）才是第 3 章证明收缩与翻转时真正使用的版本：绝对形式要求 $X$ 射影，而收缩的底空间只是一般簇。
\end{enumerate}

\subsection*{三、与我原有理解的出入}

\begin{itemize}
\item 我原以为“终端奇点”是技术上不得不接受的妥协，读完 \S2.1 才明白它是一个\textbf{刻意的选择}：终端奇点恰好是“使 $K$ 仍为 $\Q$-Cartier、且翻转不再引入更坏奇点”的最小类。换句话说，奇点类不是被发现的，而是被\textbf{筛选}出来的。
\item 我原以为锥定理的陈述（第 3 章）与收缩的存在性是两件事。本章的四步流程说明两者必须同时成立才能称“纲领可执行”；而终止性（三维由森 1988 年解决，高维至今开放）与翻转存在性是仅存的两个缺口，本书 3.45、3.46 两段也诚实地标了出来。
\item 我原以为“消没定理”是黑箱。\S2.4 的证明路线——Hodge 理论把全纯上同调抬到拓扑上同调，再用 GAGA 与循环覆盖化归到 Serre 消没——让我第一次看清：小平消没的实质是“用拓扑压住全纯”，而不是某种解析技巧。
\end{itemize}

\subsection*{四、未解疑问}

\begin{enumerate}
\item 例 2.1 的三维商奇点 $\frac1r(1,a,r-a)$ 在什么条件下不属于终端类？书中只给了“一般不终端”的例子，没有给出判据的完整证明，我打算在第 5 章 \S5.3 回头核对。
\item 命题 2.42 与定理 2.44 证明 dlt 的稳定性，用的是消解上的组合论证；但“翻转之后 dlt 保持”是否可以不依赖消解、直接从差异系数的赋值定义得到？本书没有这样处理。
\item Kawamata--Viehweg 消没定理的系数条件中，$\sum D_i$ 必须是 snc 除子；命题 2.71 说曲面情形可以去掉 snc。三维以上是否有类似的放宽？书中未涉及，值得记为后续阅读线索。
\end{enumerate}

\section*{研读札记 03 \textperiodcentered{} 第 3 章 锥定理（原书 p.95--138）}
\addcontentsline{toc}{section}{研读札记 03 \textperiodcentered{} 第 3 章 锥定理（原书 p.95--138）}
\markboth{研读札记 03 \textperiodcentered{} 第 3 章 锥定理（原书 p.95--138）}{}

\begin{quote}\small 这是全书的技术心脏。第 1 章的锥定理只对非奇异簇成立，且证明给不出"极射线可收缩"；第 3 章用完全不同的方法（无基点方法）把锥定理做到 klt 对，并顺带得到收缩定理、相对版本与 MMP 的精确表述。以下是我按节顺序做的笔记。\end{quote}

\subsection*{一、先说清楚：这一章到底在证什么}

作者在 \S3.1 一次列出四个定理，它们的关系不是并列，而是\textbf{一条推导链}：

\[
\text{非消没定理（3.4）}\Rightarrow\text{无基点定理（3.3）}\Rightarrow\text{有理性定理（3.5）}\Rightarrow\text{锥定理（3.7）}
\]

但他\textbf{先证无基点定理}（\S3.2），因为"这个证明最清楚地表达了思想"。这种"按理解难度而非逻辑依赖排课"的写法，我读下来是本章最值得学习的地方：先给读者一个能握住的锤子，再告诉他这把锤子从哪来。

\begin{itemize}
\item 定理 3.1 = 第 2 章消没定理 2.64 的改写：$Y$ 非奇异、$D=L+\sum d_iD_i$ nef 且大、$\sum D_i$ 是 snc $\Rightarrow$ $H^i(Y,\mathcal{O}_Y(K_Y+\lceil D\rceil))=0$（$i>0$）。注意 $\lceil D\rceil$ 是逐项取上整，\textbf{$D$ 本身不是整系数除子，不能直接代入}——这是审稿时最容易出错的地方。
\item 定理 3.3（无基点）：$(X,\Delta)$ 固有 klt、$\Delta$ 有效、$D$ nef Cartier，且某个 $a>0$ 使 $aD-K_X-\Delta$ nef 且大 $\Rightarrow$ 对 $b\gg0$，$|bD|$ 无基点。\textbf{"丰富 + 典范"的配方}从此贯穿全书。
\item 定理 3.4（非消没）：$X$ 固有、$D$ nef Cartier、$G$ 为 $\mathbb{Q}$-除子，若 $aD+G-K_X$ nef 且大、$(X,-G)$ klt $\Rightarrow$ 对 $m\gg0$，$H^0(X,mD+\lceil G\rceil)\ne0$。它回答"什么时候一定有截面"，是链的源头。
\item 定理 3.5（有理性）：临界斜率 $r(H)=\max\{t:H+t(K_X+\Delta)\ \text{nef}\}$ \textbf{是有理数，且分母被 $a(X)(\dim X+1)$ 界住}。
\item 定理 3.7（锥定理）：$\overline{NE}(X)=\overline{NE}(X)_{(K_X+\Delta)\ge0}+\sum\mathbb{R}_{\ge0}[C_j]$，$0<-(K_X+\Delta)\cdot C_j\le2\dim X$，且极锥面可收缩、线丛可下降。
\item 补充 3.6：$r(H)$ 达到临界时，\textbf{恰有一条 $(K_X+\Delta)$-负极射线}与之正交。这句话是把有理性定理变成锥定理的"转轴"。
\end{itemize}

\subsection*{二、无基点方法：一个可以背下来的三段式}

\S3.9 的"基本方针"我反复读了三遍，觉得它是全章最值钱的一页。思想极简单：

\begin{enumerate}
\item \textbf{想从 $X$ 上的除子拿到 $F$ 上的截面}：用 $0\to\mathcal{O}_X(K_X+M)\to\mathcal{O}_X(K_X+M+F)\to\mathcal{O}_F(K_F+M|_F)\to0$，若 $K_X+M$ 满足消没，就得到满射 $H^0(X,\mathcal{O}(K_X+M+F))\twoheadrightarrow H^0(F,\dots)$。
\item \textbf{把目标凑成上面的形状}：要证 $|bD|$ 无基点，就想把 $bD-K_X$ 写成 $M+F$（$M$ nef 且大，$F$ 不可约），这样 $bD|_F-K_F=M|_F$，满射正是"$F$ 上有截面 $\Rightarrow$ $X$ 上有截面"。
\item \textbf{凑不出来就爆破}：改为在 $f:Y\to X$ 上做，把 $bf^*D-K_Y$ 写成 $(\text{nef 且大})+F-A$，其中 $\lceil A\rceil$ 是有效的 $f$-例外除子，于是 $H^0(X,bD)=H^0(Y,bf^*D+\lceil A\rceil)$，\textbf{例外部分自动被"吸收"}。
\end{enumerate}

真正的手艺在第 3 步之后：写 $f^*|mD|=|L|+\sum r_jF_j$（移动部分 + 固定部分），取 $c>0$ 使

\[
\min_j(a_j-cr_j)=-1,
\]

再\textbf{用系数极小的 $\sum p_jF_j$ 把所有东西扰动一下}，使取到最小值的 $j$ 唯一——于是 $F$ 恰好一个不可约除子，而 $A$ 不含 $F$ 为其成分。定义

\[
N(b,c):=bf^*D-K_Y+\sum(-cr_j+a_j-p_j)F_j
\]

它被拆成三项：$(b-cm-a)f^*D$（nef）、$c(mf^*D-\sum r_jF_j)$（无基点）、$f^*(aD-K_X-\Delta)-\sum p_jF_j$（丰富）。三者之和丰富 $\Rightarrow$ 消没 $\Rightarrow$ 满射到 $F$ $\Rightarrow$ 再用非消没定理在 $F$ 上取到非零截面 $\Rightarrow$ $F\not\subset\mathrm{Bs}|bD|$ $\Rightarrow$ 与 $|mD|$ 的基点集相矛盾。

\textbf{我记住的是这个"扰动"技巧}：一个本该"可能有好几个"的极小值集合，用一次无穷小扰动变成"恰好一个"。它不改变几何，只改变账本。Kollár--Mori 反复用这一手（有理性定理的步骤 4 又用了一次），值得当工具收集。

\subsection*{三、有理性定理：几何里的丢番图逼近}

\S3.4 的证明结构很值得欣赏：\textbf{它把"$r$ 是无理数"翻译成"带状区域里落入太多格点"}。

\begin{itemize}
\item 消没定理 $H^1(X,\mathcal{O}(pH+qK_X))=0$ 成立的条件恰是 $pH+(q-1)K_X$ nef 且大，即
\end{itemize}

\[
\frac{q-1}{p}<r<\frac{q}{p}\qquad(3.1)
\]

于是在这条宽度约 $1/p$ 的带里，$\chi$ 与 $h^0$ 都非零、线性系都非空——\textbf{只要带里的点足够多，就能挑出一族基点集严格递减的线性系，与 Noether 归纳法矛盾}。
\begin{itemize}
\item 引理 3.19 是纯数论：非零多项式 $P(x,y)$（次数 $\le n$）在 $0<ay-rx<\epsilon$ 的所有充分大整数解上取零 $\Rightarrow$ $r$ 有理且分母 $\le a(n+1)/\epsilon$。证明用"$P$ 与 $(y'x-x'y)$ 有 $n+1$ 个公共解 $\Rightarrow$ 后者整除 $P$"，再用更小的 $\epsilon$ 重复，$(n+1)$ 次后矛盾——\textbf{次数有限，所以只能"挡住"有限多个方向}。
\item 多项式从哪来？$P(u_1,\dots,u_k)=\chi(\sum u_iD_i+\lceil A\rceil)$，是 Hilbert 多项式；由消没 + 非消没可知它在某些点上非零，故不恒为零。
\end{itemize}

我在这里学到的思维方式是：\textbf{几何的"临界值"如果无理性，就会强迫离散对象（格点）密度过高，从而与有限次数的代数对象（多项式）冲突。} 这是一种"刚性"论证，和 Roth 定理的精神相通，只是这里用初等办法就够。

\subsection*{四、锥定理：纯凸几何 + 有理性定理}

\S3.3 的思路（3.14）很短：对丰富除子 $H$，$m(H+r(H)(K_X+\Delta))$ 是 nef 但\textbf{不丰富}的，于是 $\overline{NE}(X)\cap(H+r(K_X+\Delta))^\perp$ 是曲线锥的一个真面；变动 $H$ 就得到各种面。形式化后就是定理 3.15（抽象凸几何引理）。

证明的四步是标准的"降维打击"：
\begin{enumerate}
\item 对 nef 的 $L$，令 $F_L=L^\perp\cap\overline{NE}$；若 $F_L\not\subset\overline{NE}_{K\ge0}$，构造 $D(nL,H)=na(K)L+a(K)H+r_L(H)K$，得非平凡面 $F_{D(nL,H)}\subset F_L$。
\item 选 $H_i$ 使它们在 $F_L$ 上的限制构成对偶基，可令 $\dim F_{D(nL,H_i)}<\dim F_L$ ——\textbf{只要面的维数 $>1$，就能再切一刀}。
\item 迭代到 $\dim F_L=1$，证明 $\overline{NE}$ 由 $K\ge0$ 部分与这些 1 维面共同生成（否则用分离超平面造出矛盾，手法同定理 1.24）。
\item \textbf{离散性}：在仿射坐标 $\xi\mapsto\big((\xi\cdot H_1)/(\xi\cdot K),\dots\big)$ 下，这些 1 维面上的点坐标分母被 $a(K)$ 界住，故不聚积；于是对任意 $\epsilon>0$，$F_L\cdot(K+\epsilon H)<0$ 的 $F_L$ 只有有限个，右侧是闭锥。
\end{enumerate}

第 7、8 步接上无基点定理：极锥面的支撑除子 $D$ 取 $m\gg1$ 使 $mD-(K_X+\Delta)$ 丰富 $\Rightarrow$ $|mD|$ 无基点 $\Rightarrow$ Stein 分解给出\textbf{唯一}的收缩态射 $\mathrm{cont}_F$；再证明线丛可下降（定理 3.7(4)）。

\subsection*{五、MMP 的执行与"终点"}

\S3.7 把前面所有东西组装成算法，并给出我一直想确认的几条"账本"：

\begin{itemize}
\item 收缩一条极射线：$\Pic$ 正合列 $0\to\Pic Z\to\Pic X\to\mathbb{Z}\to0$，故 $\rho(Z)=\rho(X)-1$（推论 3.17）。\textbf{每一步恰好掉一个 Picard 数}——这是终止性的"计数器"。
\item 除子收缩与 Fano 收缩都保持 $\mathbb{Q}$-分解性（推论 3.18）；翻转也保持 $\mathbb{Q}$-分解性与奇点类别（推论 3.42--3.44 一带）。
\item \textbf{3.45（终止性）与 3.46（翻转存在性）是纲领的两个"未证前提"}：终止性至今一般维数未知；翻转存在性：二维平凡（无翻转）、三维终端 [Mor88]、三维 klt [Sho92]、三维 lc [K+92]。
\item 定理 3.47（二维完全解决）：二维 $\mathbb{Q}$-分解 dlt 对经至多 $\rho(X)-1$ 次收缩，终点只有三种——对数极小模型、对数直纹曲面、对数 Del Pezzo 曲面。\textbf{二维没有翻转，所以二维 MMP 是定理而不是纲领。}
\end{itemize}

\subsection*{六、极小模型与标准模型：定义为什么这么绕}

\S3.8 是我读得最吃力的地方，读完才明白作者为何要先举反例 3.49：设 $S$ 是只有一个奇点的曲面、$f:S'\to S$ 是极小消解、例外曲线 $(D^2)=-n$，则 $K_{S'}\equiv f^*K_S-\frac{n-2}{n}D$。对 $(S',cD)$ 跑 MMP：$c>\frac{n-2}{n}$ 时第一步就收缩 $D$，$c\le\frac{n-2}{n}$ 时不收缩。\textbf{同一个底空间 $S$，可以是 $(S,0)$ 的极小模型，却未必是 $(S',cD)$ 的极小模型}——差别系数 $c$ 改变了"极小"的含义。

所以定义 3.50 必须比较差异系数，并列出五条：$f^w$ 固有、$\varphi^{-1}$ 无例外除子、$\Delta^w=\varphi_*\Delta$、$K_{X^w}+\Delta^w$ 为 $\mathbb{Q}$-Cartier 且 $f^w$-nef、以及\textbf{对每个 $\varphi$-例外除子 $E$：$a(E,X,\Delta)\le a(E,X^w,\Delta^w)$}。再加"严格小于"就是伪极小模型，加"$f$-丰富"就是伪标准模型；lc 有效时叫（弱）标准模型，dlt 时叫极小模型。

三条结论：
\begin{itemize}
\item \textbf{命题 3.51}：把 (3) 换成更弱的 $(3')$（$\varphi_*\Delta-\Delta^w$ 去掉 $\varphi^{-1}$-例外后有效），仍得到 $(5')$——所以弱伪标准模型"不会让奇点变坏"。证明是标准除子比较 $B:=g^*(K_X+\Delta)-h^*(K_{X^w}+\Delta^w)\equiv E_2-E_1$，$-B$ 为 $g$-nef、$g_*B$ 有效，用引理 3.39 得 $E_2-E_1$ 有效。
\item \textbf{定理 3.52}：伪标准模型唯一，且 $X^c=\operatorname{Proj}_S\bigoplus_{m\ge0}f_*\mathcal{O}_X(mK_X+\lfloor m\Delta\rfloor)$；两个伪极小模型在余维数 1 处同型。\textbf{"唯一性在余维数 1 处成立"是 MMP 的标准口径}，不能指望更强（翻转正是余维数 2 的现象）。
\item 注意环里用 $\lfloor m\Delta\rfloor$ 而不用 $\lceil m\Delta\rceil$：因为 $\lfloor m_1\Delta\rfloor+\lfloor m_2\Delta\rfloor\le\lfloor(m_1+m_2)\Delta\rfloor$ 才有环结构（次可加性），换成上整不等式反向，连环都不是。这个细节我在译定理 3.11 时专门加了批注。
\end{itemize}

\subsection*{七、我的疑问 / 想回头查的}

\begin{enumerate}
\item 引理 3.19 的证明中"重复 $(n+1)$ 次即得矛盾"这一步，我按逻辑补全了（每次用更小的 $\epsilon$，得到不同的整除因子，$(n+1)$ 个因子相乘次数超过 $n$），但原书写得很简；我打算对照英文原著确认没有漏掉"$\epsilon$ 需重新选取"的条件。
\item 定理 3.15 步骤 1 里 $r_L(n,H)$ 是"非负整数"——靠有理性定理给分母 $a(K)$，再乘 $a(K)$ 得整数；但 $n\to\infty$ 时它收敛到 $r_L(H)$ 的论证依赖 $F_L\setminus\overline{NE}_{K\ge0}\ne\varnothing$ 的选取，我补了一条批注说明。\textbf{这类"书里一步跨过、读者要自己补"的地方，正是译注最该发光的地方。}
\item \S3.7 与 \S3.8 的次序：先用锥定理跑 MMP（3.31），再用 3.50 定义"MMP 的终点"，逻辑上是"先跑再说终点长什么样"。初读容易倒过来。
\item 本章反复出现的 $a(X)$（使 $a(X)(K_X+\Delta)$ 为 Cartier 的最小正整数）与 \S3.6 相对版本里的 $a(K)$ 是同一角色；书中不总写全，读时要把"对某个固定整数取分母"这句话记在心上。
\end{enumerate}

\subsection*{八、如果只让我记一件事}

\textbf{锥定理的本质是一句凸几何 + 一句数论：曲线锥的负部由可数条极射线生成，而极射线在负半空间里是离散的；离散性的来源是"临界斜率的分母有界"。} 有了它，MMP 才是一个"每一步都知道往哪走、且计数器在减少"的算法，而不是一堆技巧的堆叠。

\section*{研读札记 04 \textperiodcentered{} 第 4 章 极小模型纲领的二维奇点（原书 p.139--190）}
\addcontentsline{toc}{section}{研读札记 04 \textperiodcentered{} 第 4 章 极小模型纲领的二维奇点（原书 p.139--190）}
\markboth{研读札记 04 \textperiodcentered{} 第 4 章 极小模型纲领的二维奇点（原书 p.139--190）}{}

这一章是全书唯一的“完全分类”章节，也是我读得最踏实的一章：所有结论都可以在纸上画图验证。它表面上讲二维，实际是为第 5、6、7 章提供局部模型——三维的翻转最终都归约为“二维 Du Val 奇点 + 一个离散不变量”。

\subsection*{一、本章在做什么}

五节各有独立目标，但共用一套语言（极小消解 + 差异系数 + 例外图的交点矩阵）：

\begin{itemize}
\item \S4.1 用极小消解把二维 lc 对 $(X,\Delta)$ 的局部结构一次算清：$X$ 非奇异、或为循环商奇点、或为一条链上的商奇点，$\Delta$ 的系数在 snc 部分之外只能是 $1-\frac1n$ 这类离散值。
\item \S4.2 把 $K_X$ Cartier 的二维标准奇点等同于 Du Val 奇点，给出 $A_n,D_n,E_6,E_7,E_8$ 的方程与例外图。
\item \S4.3 证明 Du Val 奇点的极小消解可以\textbf{同时}进行：存在一个有限态射 $X'\to X$，使 $X'\times_X Y\to X'$ 光滑（定理 4.28 的局部形式，定理 4.43 的整体形式）。
\item \S4.4 处理二维椭圆奇点，用例外循环 $Z$ 把问题压到曲线上，得到三条标准方程。
\item \S4.5 讨论孤立超曲面奇点的半普遍形变：用 Artin 逼近与 Tougeron 引理把方程写成 $f+\sum t_ig_i$。
\end{itemize}

\subsection*{二、我抓到的技术要点}

\begin{enumerate}
\item \textbf{差异系数在二维可以“看图算出”。} 设 $f\colon Y\to X$ 为极小消解，$K_Y=f^*K_X+\sum a_iE_i$，则 $a_i=-1-\dfrac{(E_i\cdot E_i)+2}{2}$（对 $X$ 非奇异的爆破链逐次计算）。Du Val 的判据 $K_X$ Cartier 与 $a_i=0$ 等价，于是 $A$--$D$--$E$ 图恰好是“每个顶点自交 $-2$”的负定树。这一条我在读第 2 章时只觉得是定义，这里才看到它是\textbf{可计算}的。
\item \textbf{$k=-(Z\cdot Z)$ 是椭圆奇点唯一的离散不变量。} 记号 4.44 取 $Z=\sum\lceil-a_i\rceil E_i$（注意是上取整！），则 $Z$ 数值 $f$-平凡，$k=-(Z\cdot Z)\ge1$；$k\ge3$ 时 $X$ 是 $\mathbb{A}^k$ 中的超曲面，$k=2$ 时是 $z^2+q_4(x,y)=0$，$k=1$ 时是 $z^2+y^3+yq_4(x)+q_6(x)=0$。用上取整而非下取整，正是为了让 $Z$ 的系数为整数：若是 $\lfloor-a_i\rfloor$，$Z$ 就不整了。这一点与我读第 3 章时“环要用下取整 $\lfloor m\Delta\rfloor$”的印象相反，值得警惕——\textbf{同一个 $K_Y$ 分解，取整方向由用途决定}。
\item \textbf{同时消解的本质是“基变换”。} 定理 4.28 的证明不是逐纤维消解，而是构造一个有限覆盖 $X'\to X$，使拉回族的消解成为整体的光滑态射。这个思想在第 7 章半稳定归约里会以“半稳定基变换”的形式重演，我读到这里先记下：\textbf{先改基，再消解}。
\item \textbf{形变理论在本章只作工具。} \S4.5 的半普遍形变空间存在性（Grauert）与显式方程 $f+\sum t_ig_i$（定理 4.61）本身不难，难的是 Tougeron 引理在解析、形式、代数三种设定下的差别：解析情形用 Artin 逼近，代数情形要额外假定。本书的处理次序（先解析、再形式、后代数）是标准做法，值得照抄。
\end{enumerate}

\subsection*{三、与我原有理解的出入}

\begin{itemize}
\item 我一直以为椭圆奇点比 Du Val 奇点“更坏”，所以在第 5 章里应该被排除。实际上定理 5.35 恰恰用椭圆奇点\textbf{描述}三维标准奇点：三维标准奇点的横截切片就是 Du Val 或椭圆奇点。分类越彻底，越能当字典用——这是本章方法论上给我的最大教益。
\item 我以为“同时消解”是为了让翻转的构造更方便。真正的用途是\textbf{形变}：Du Val 奇点族的纤维要同时光滑化，才能保证第 7 章半稳定翻转中的族在每一步都几何整。
\item 关于 $Z$ 的两种定义（记号 4.44 的 $\sum\lceil-a_i\rceil E_i$ 与定理 4.57 的 $\omega_Y=f^*\omega_X(-Z)$），我起初以为是笔误。核对英文原著后才确认：$K_X$ 非 Cartier 时后者只是数值 $f$-平凡的例外除子，二者的差别正是“$\Q$-Cartier 与 Cartier”的差别。这一点本书未加批注，我已在译文中加了 Tutor Note。
\end{itemize}

\subsection*{四、未解疑问}

\begin{enumerate}
\item 定理 4.22（$A$--$D$--$E$ 分类）的证明被引理 4.23 打断，我按译者建议分块闭合了。但分类的“唯一性”部分（同一个奇点不能同时属于两个类型）书中只用图形论证，是否可以从交点矩阵的正定性直接推出？
\item 椭圆奇点在 $k=1$ 时的标准方程含 $q_4,q_6$，其模空间维数（Tjurina 数 $\tau$）与 $k$ 的关系书中未讨论，\S4.5 只在注里提了一句。
\item 同时消解在特征 $p$ 下是否成立？本书全程在 $\mathbb{C}$ 上，但第 1 章的折弯引理用过特征 $p$，两处的“特征”用法完全不同，容易混淆。
\end{enumerate}

\section*{研读札记 05 \textperiodcentered{} 第 5 章 极小模型纲领的奇点（原书 p.191--232）}
\addcontentsline{toc}{section}{研读札记 05 \textperiodcentered{} 第 5 章 极小模型纲领的奇点（原书 p.191--232）}
\markboth{研读札记 05 \textperiodcentered{} 第 5 章 极小模型纲领的奇点（原书 p.191--232）}{}

这一章是全书的“定义装配车间”：不发明新方法，只挑出恰好够用的奇点类，再证明它们在收缩、翻转、形变之下都不跑出去。读完最大的感受是：\textbf{MMP 的可执行性，本质上等于“奇点类在运算下封闭”这条纯技术事实}。第 2 章给了纲领的骨架，第 6、7 章给了翻转的构造，而这两者的合法性全部抵押在本章。

\subsection*{一、本章在做什么}

\begin{itemize}
\item \S5.1 有理奇点：给出三条等价刻画，并建立迹与范数（trace/norm）的延拓工具。商奇点是最典型的例子。
\item \S5.2 dlt 对的局部理论，主定理 5.22：dlt 对的有理性。由此可以把 $\Q$-Cartier 的 Weil 除子当 Cartier 用（命题 5.26、推论 5.27）——这是全书所有“在 dlt 对上做除法”的论证的许可证。
\item \S5.3 三维终端奇点与标准奇点：用二维 Du Val 奇点、椭圆奇点作为横截切片来分类（定理 5.34、5.35），并给出 cDV 与循环商的描述；表 5.1 列出终端奇点的解析型。
\item \S5.4 伴随与伴随反转：定理 5.50（inversion of adjunction）把 $(S,B|_S)$ 的 klt 性等价于 $(X,S+B)$ 的 plt 性，是归纳到低维的核心工具。
\item \S5.5 CM 层对偶理论：不进入导出范畴，只用一个层的 $\mathcal{H}om$ 建 Serre 对偶与伴随公式。
\end{itemize}

\subsection*{二、我抓到的技术要点}

\begin{enumerate}
\item \textbf{有理奇点的“三等价”把上同调问题变成结构层问题。} 定理 5.10：$X$ 有理奇点 $\Longleftrightarrow$ $\cO_X\to Rf_*\cO_Y$ 同型（对某个/任一消解）$\Longleftrightarrow$ $Y$ CM 且 $f_*\omega_X\cong\omega_Y$ 且高次消没。第三条最有用：它把“奇点坏不坏”翻译成“对偶化层有没有掉”。我读第 2 章时一直把“有理”当成一个孤立定义，到这里才看出它其实是\textbf{判据}。
\item \textbf{迹与范数靠 $S_2$ 延拓。} 定义 5.6 里，先在余维 $\ge2$ 的 $Z$ 之外定义 $\operatorname{Trace}$（那里层局部自由），再用 Hartogs 现象延拓到全空间。前提是 $X$ 为 $S_2$（正规即 $S_2+R_1$）。这解释了为什么本章开头要花两页讲 CM 与 $S_d$：不是为了堆砌定义，而是\textbf{延拓定理的门槛}。
\item \textbf{dlt 的有理性几乎是“免费”的，klt 更是。} 定理 5.22 的证明走归纳：沿边界成分的伴随、再对维数归纳。它之所以重要，是因为在 dlt 对里每一步翻转都会产生新的边界，而这些边界必须仍是 dlt，否则下一步就没法用消没定理。
\item \textbf{三维终端奇点 = cDV 奇点 / 有限群。} 定理 5.34、5.35 用“横截切片是 Du Val 或椭圆奇点”来给终端/标准奇点分类，表 5.1 的 cA、cAx、cD、cE 是\textbf{上同调型}的记号（$H^i=\frac1r(1,a,r-a)$ 之类的循环商），不是奇点的终名。作者明说本书不需要这个分类，只把它当背景——这一点很值得学：\textbf{分类做得再细，也要敢于声明“本书不用”}。
\item \textbf{伴随 vs 伴随反转：方向相反。} 命题 5.46 是伴随方向的不等式（从 $X$ 推到 $S$，$a(E,S,B|_S)\ge a(E,X,S+B)$），定理 5.50 才是反转（从 $S$ 反推 $X$）。我第一遍读时把两者记混，直到看清证明用的是“对 $S$ 上的每个赋值构造 $X$ 上的例外除子”才分开。
\item \textbf{对偶理论只用一个层。} \S5.5 的写法刻意避开导出范畴：Serre 对偶用 $\mathcal{H}om$ 与 CM 层直接陈述（定理 5.71），伴随公式写成 $\omega_X\cong\omega_Y\otimes\det(\mathcal{I}/\mathcal{I}^2)^\vee$ 这类显式形式。对读者（比如我）而言这比 $\mathbf{R}\mathcal{H}om$ 友好得多，代价是证明里要反复验证 CM 条件。
\end{enumerate}

\subsection*{三、与我原有理解的出入}

\begin{itemize}
\item 我原以为“终端奇点”的定义只关于差异系数，读 \S5.3 才发现三维情形还要求\textbf{指数有界}（表 5.1 的 $r$ 有限），这一步靠的是 Mori 的终端奇点有界性。分类的有限性不是定义的推论，而是\textbf{定理}。
\item 我原以为伴随反转只是一个不等式技巧。定理 5.50 的等号情形说明它其实是\textbf{等价}：$(S,B|_S)$ klt 当且仅当 $(X,S+B)$ plt。等价的用处在于“反向”使用——先在一个低维切片上验证 klt，再断言整个对是 plt，这正是第 7 章半稳定翻转归纳的骨架。
\item 我原以为 CM 只是上同调的技术条件。定义 5.1 的 (1)(2) 等价（正则列 vs 维数条件）说明 CM 是“维数不塌陷”的代数表达，而 $S_2$ 正是延拓定理所要求的最小“不塌陷”程度。
\end{itemize}

\subsection*{四、未解疑问}

\begin{enumerate}
\item 定理 5.22 的证明对射影簇写得很完整，对一般代数/复解析情形只引文献。这个“局部陈述、射影证明”的模式中，真正的障碍是相对对偶理论的表述难度——书中第 5 章开头也承认了。我在译注里加了提醒，但自己仍未完全消化相对对偶的构造。
\item 表 5.1 的编号（cA/n、cAx/4、cD/n、cE/n）与定理 5.34 的对应关系，书中只给了一张表，没有逐项列出方程。若要做“查表”式的阅读，还须对照 [Mor85] 或 [Rei87a]。
\item 定理 5.34 与 5.35 的证明集中在编号段 5.36，两者的陈述末尾印 $\square$（表示证明另给），这与“书中不证”（$\blacksquare$）容易混淆。这一点我在页图上逐条核对过，译文按页图保留。
\end{enumerate}

\section*{研读札记 06 \textperiodcentered{} 第 6 章 三维翻模（原书 p.233--256）}
\addcontentsline{toc}{section}{研读札记 06 \textperiodcentered{} 第 6 章 三维翻模（原书 p.233--256）}
\markboth{研读札记 06 \textperiodcentered{} 第 6 章 三维翻模（原书 p.233--256）}{}

这一章札记记录我读第 6 章时的理解路径。第 6 章是全书唯一一章"专门解剖一个技术装置"的：\S6.1 讲翻转与翻模的一般理论（唯一性、局部性、与分次代数的关系），\S6.2 证三维终端翻模的存在性与终止性，\S6.3 讲终端化与 Q-分解，\S6.4 用"crepant 降下"把标准翻模归约到终端翻模。

\subsection*{一、这一章到底在补什么缺口}

MMP 的每一步只有三种动作：\textbf{除子收缩}（余维 1）、\textbf{纤维型收缩}（相对维数 1）与\textbf{小收缩}。前两种收缩的终点仍然是"正常的"簇，第三种不是：$K_X$ 在小收缩 $f:X\to Y$ 上为负时，$Y$ 上 $K_Y$ 甚至不是 Q-Cartier，MMP 在这里断掉。翻转 $f^+:X^+\to Y$ 就是把断点接上，$X^+$ 上 $K_{X^+}$ 是 $f^+$-丰富的。所以第 6 章的整个存在性就是"MMP 能不能一直走下去"。

我此前的一个误解是：以为翻转是"把负曲线切掉换一条"。不对。翻转的核心是\textbf{同一基 $Y$ 上换一个模型}，$X\dashrightarrow X^+$ 在 $Y$ 上同构于 $Y$ 的同一开集，只是例外集上的纤维被换掉。它不切除任何东西，几何上的"实心"内容全部藏在 $Y$ 的非 Q-Cartier 性里。

\subsection*{二、翻转的存在性 = 一个分次代数的有限生成性（\S6.1）}

引理 6.2（[Kaw88]）是本章的引擎：对 $Y$ 上的 Weil 除子 $B$，$R(Y,B)=\bigoplus_{m\ge0}\cO_Y(mB)$ 有限生成 $\iff$ 存在 $g:Z\to Y$ 使 $\operatorname{Ex}(g)$ 余维 $\ge2$、$B'=g_*^{-1}B$ 是 Q-Cartier 且 $g$-丰富。推论 6.4 把它特化到翻转：\textbf{$(K+D)$-翻转存在 $\iff$ $\bigoplus_{m\ge0}\cO_Y(mf_*(rK_X+rD))$ 有限生成}。

这里有两个我当时想了一会儿的点：

\begin{enumerate}
\item 引理 6.2 的证明里，从 (2) 推 (1) 那一步用的是"$\operatorname{Ex}(g)$ 余维 $\ge2$ 故截面可沿 $g$ 延拓"——这正是代数几何里所有"余维 2 处无碍"论证的同一手法（与 $S_2$ 延拓、Hartogs 一个道理）。它把 $g_*\cO_Z(mB')\to\cO_Y(mB)$ 变成同构，然后 $Z=\Proj_Y R$ 的唯一性随之而来。
\item 从 (1) 推 (2) 那一步是反证：若 $g$ 有例外除子 $E$，则 $\cO_Z\subseteq\cO_Z(E)$ 给出单射 $\tau:\cO_Y(mB)\to f_*(\cO_Z(m)(E))$，在 $g(\operatorname{Ex}(g))$ 之外是同构，但 $\cO_Y(mB)$ 自反（余维 2 集之外同构就够）——矛盾。\textbf{"自反"在这里的角色就是"余维 2 处不能有伪装的例外"}。
\end{enumerate}

推论 6.4 还顺手给了翻模的三条结构性事实：翻转与拟标准模型一致（(1)）、唯一（(3)）、在 $\sim_f$ 意义下不依赖 $D$（(4)）。第 (4) 条是后面翻模定义的根据。

\subsection*{三、翻转是局部问题（\S6.1 后半）}

命题 6.6 说：$R(Y,B)$ 在某点的\textbf{完备化}上有限生成，等价于在该点的某个 Zariski 邻域上有限生成。推论 6.7 由此得：翻转存在只需逐点验证。这是整个存在性理论里最重要的一条"降维"：三维终端翻模的存在性只需要在 $(Q\in Y)$ 的解析芽上证明——于是三维终端奇点的分类（推论 5.39，即 $(\C^4,0)$ 的超曲面 $u^2=f(x,y,z)$）才能插手进来。

命题 6.8（[Kaw88, 3.2]）是有限覆盖下的不变性：$h:Y'\to Y$ 有限满射时，$R(Y,B)$ 与 $R(Y',h^*B)$ 的有限生成性互相等价。证明里"取 Galois 闭包再取 $G$-不变量"那一步，把有限态射的问题化为群作用的不变量问题，是代数几何里处理"指数"的标准套路。

\subsection*{四、翻模：定义 6.10 的几个微妙处}

翻模收缩 $f:X\to Y$ 是固有双有理且 $\operatorname{Ex}(f)$ 余维 $\ge2$、$K_X\equiv_f 0$。这是\textbf{小收缩里标准类为平凡}的那一类。翻转与翻模的差别就一个符号：翻转 $-(K_X+D)$ 为 $f$-丰富（$K$ 沿收缩为负），翻模 $K_X\equiv_f0$。

定义 6.10 里有一段容易被读漏的话：若 $(X,0)$ 是终端、$D$ 是 Q-Cartier 的，则 $0<\varepsilon\ll1$ 时 $(X,\varepsilon D)$ 仍终端，于是\textbf{终端翻模可以看作终端翻转的特例}；但即使 $(X,0)$ 是标准的，也可能对任何 $\varepsilon$ 都有 $(X,\varepsilon D)$ 不标准——所以\textbf{标准翻模不是标准翻转的特例}。这句话是定理 6.17（终止性）处理的形态差异的来源：标准翻模只能单独证终止。

\subsection*{五、三维终端翻模的存在性（定理 6.14）}

证明路线我梳理成三步：

\begin{enumerate}
\item \textbf{先证基 $Y$ 自身终端}：$f$ 是 $\overline{NE}(X/Y)$ 中一个 $(K_X+\varepsilon D)$-负极端面的收缩，由定理 3.25(4) 得 $K_Y$ Q-Cartier 且 $K_X=f^*K_Y$，再用引理 2.30 得 $(Y,0)$ 终端。
\item \textbf{归约到解析芽}：由推论 6.7 只须证 $R(Y,K_Y+f_*D)$ 在 $(Q\in Y)$ 上有限生成。
\item \textbf{用分类把芽变成超曲面}：由推论 5.39，$(Q\in Y)$ 有有限覆盖 $\pi:(\tilde Q\in\tilde Y)\to(Q\in Y)$ 使 $\tilde Y$ 是 $u^2=f(x,y,z)$。此时 $\tilde Y$ 还是 $(\C^3,0)$ 的二重覆盖，带 2 阶自同构 $\iota$。关键观察：对 $\tilde Y$ 上 Weil 除子 $F$ 有 $\iota^*F\sim-F$（因 $F+\iota^*F=g^*g_*F\sim0$），于是 $R(\tilde Y,F)\cong\iota^*R(\tilde Y,-F)$。又 $R(Y,-K_Y-f_*D)$ 有限生成（$-K_X-D$ 为 $f$-丰富），沿 $g$ 拉回、沿 $\iota$ 翻符号，得到 $R(\tilde Y,\pi^*f_*D)$ 有限生成，再沿 $\pi$ 下降得 $R(Y,K_Y+f_*D)$ 有限生成。
\end{enumerate}

我最欣赏第 3 步里的\textbf{符号翻转}：把"$K$ 沿收缩为负"换成"$K$ 沿二重覆盖平凡"的对称性，一举把正性换到另一边。这也是为什么四维以上不知情：四维终端奇点没有这样可计算的分类，$\iota$ 这个自同构就没有来源。

定理 6.15（翻模保持解析奇点）是 6.14 的副产品，但对后续极其重要：奇点类型在翻模下不变，所以标准翻模的构造可以逐点归约到局部模型上——这是终止性证明能"逐点记账"的前提。

\subsection*{六、终止性（定理 6.17）}

不变量是定义 6.20 的 $d(X,\Delta)$：对标准对 $(X,\Delta)$，令 $a=\max a_i$、$S=\sum a_i\Z_{\ge0}$，
\textbackslash{}[
d(X,\textbackslash{}Delta)=\textbackslash{}sum\_\{\textbackslash{}xi\textbackslash{}in S,\textbackslash{} \textbackslash{}xi\textbackslash{}ge a\}\textbackslash{}\#\textbackslash{}\{E:\textbackslash{} a(E,X,\textbackslash{}Delta)<1-\textbackslash{}xi\textbackslash{}\}.
\textbackslash{}]
它是有限的（因 $\operatorname{totaldiscrep}=-a$ 与命题 2.36），且在翻转下不增（引理 3.38）。

证明的结构是\textbf{对边界除子个数 $k$ 归纳}：
\begin{itemize}
\item $k=0$：每次翻转给出 $d(X^{j-1},0)>d(X^j,0)$（引理 6.21），$d$ 是自然数，故只有有限次。
\item $k>0$：若后翻转曲线 $C^j$ 落在 $D^j_k$ 中，则 $a_k<1$ 且 $d$ 下降；故 $j\gg0$ 时 $D^j_k$ 不含后翻转曲线。于是 $\bar D^j_k$ 随后 $j$ 稳定（引理 3.40：例外曲线数值独立 $\Rightarrow$ 同构），也就最终不含前翻转曲线。剩下的翻转其实是关于 $\sum_{i<k}a_iD^j_i$ 的翻转，归到 $k-1$。
\end{itemize}

这就是"MMP 不会无限翻"的账本：\textbf{每次翻转要么消耗离散不变量 $d$，要么让某个边界分量脱离游戏}。和二维 MMP 的终止性（每次收缩让 $\rho$ 降 1）相比，三维这本账记在"赋值表"上而不是"锥的维数"上——维数在三维不再提供有限性。

定理 6.17 还有一条附加项："三维射线型标准 $(K_X+\Delta)$-翻转的任意列都有限"。注意它对\textbf{解析}空间要限制在紧集的小邻域（注意 6.18），否则翻转曲线可以跑向无穷远，序列本身就不是"有极限"的对象。

\subsection*{七、终端化与 Q-分解（\S6.3）}

\begin{itemize}
\item \textbf{终端化（定理 6.23）}：只具标准奇点的三维簇 $X$ 存在 crepant 射影双有理 $f:Y\to X$，$Y$ 至多终端。证明按 $e(X)$（$a(E,X)=0$ 的例外除子个数）归纳：若某点的指数 1 覆盖不是 cDV，由定理 5.35 有射影 $Y\to X$ 把至少一个除子收缩到 $P$，$e$ 下降；若全是 cDV 且 $\dim\Sing X=1$，由定理 6.27 收缩到曲线，$e$ 下降；若全是 cDV 且奇点孤立，由引理 6.31 已是终端。
\item \textbf{Q-分解（定理 6.25）}：终端三维簇 $X$ 存在 $f:Y\to X$，$Y$ 终端且 Q-分解，$f$ 在余维 1 处同构。它基于定理 6.32：$R(X,D)=\bigoplus\cO_X(\nu D)$ 有限生成、并存"符号爆破" $g:Y\to X$ 使 $D'=g_*^{-1}D$ Q-Cartier。
\end{itemize}

定理 6.27 证明里我最想记住的是：\textbf{符号幂 $I^{(\nu)}$ 的有限生成性}。$I^{(\nu)}$ 只在 $C$ 的泛点处与 $I^\nu$ 相同，其余点可以更大；$I^{(\nu)}/I^{\nu}$ 支集零维、长度有限，所以"局部差"是有限信息，而整体有限生成性靠"沿 $V$ 的 Galois 基变换 + $G$-不变量"搬到 $\C^3$ 的 Du Val 情形里去做。引理 6.28（$F\sim_f D$ 的有效 Cartier 除子存在）则是把所有"$f$-线性等价"的账消掉的工具。

\subsection*{八、\S6.4 标准翻模：“crepant 降下”与有限性收网}

（1）问题结构。\S6.1--6.3 处理的都是“沿标准类为负的方向”的翻转与翻模，而\textbf{标准翻模}（$K_X\equiv_f 0$ 的小收缩）没有负性可用：$(K_X+\varepsilon D)$-MMP 只有在 $D$ 是 $f$-丰富（终端模型的相对除子）时才启动得起来。于是本节的策略是\textbf{crepant 降下}：先在带丰富除子的模型之间把双有理映射分解成有限个 $D$-翻模（定理 6.38），再把标准情形当作 $D=0$ 的极限。

（2）定理 6.38 的骨架：从 $X_1$ 跑 $(K_{X_1}+\varepsilon D_1)$-MMP。由推论 3.54，$D_1$ 在余维 2 之外是丰富的，所以过程中只出现翻模；终点 $X_3$ 与 $X_2$ 在余维 1 处同构，再由引理 6.39 升格为同构。注意 $D_1=\pi_*^{-1}D_2$ 只是 $f_1$-nef（不是丰富），这恰好是引理 6.39 的第 (1) 条——\textbf{条件被设计得刚好够用}。

（3）引理 6.39 是本节唯一的“判别同构”工具，证明只用截面：取 $n\gg0$，$L_1=E_1+nD_1$ 与 $L_2=E_2+nD_2$ 都丰富（定理 1.18），而 $H^0(X_1,\cO(mL_1))=H^0(X_2,\cO(mL_2))$，故 $|mL_1|$、$|mL_2|$ 把两个模型嵌入 $Z$ 上的同一个像中，$\pi$ 只能同构。这个“用 $H^0$ 把双有理变双正则”的手法，第 3 章定理 3.52（伪标准模型的唯一性）已经用过一次；我把它记成：\textbf{收住翻转溢出靠的是 $H^0$，不是负性}。

（4）命题 6.40 是相对锥定理的边角料：设 $-D$ 为 $f$-丰富、$(X,\Delta)$ 为 klt、$K_X+\Delta\equiv_f0$，则 $NE(X/Z)$ 有限生成。证明把“$0$”拆成 $\varepsilon D+(\varepsilon/2)(-D)$，对 $K_X+\Delta+\varepsilon D+(\varepsilon/2)(-D)\equiv_f(\varepsilon/2)D$ 应用定理 3.25(2)，得 $\overline{NE}$ 由有限极射线与 $\overline{NE}_{D\ge0}$ 生成；而 $-D$ 丰富给出 $\overline{NE}_{D\ge0}=0$（定理 1.44）。这种“把平凡类拆成正负两部分、用丰富的负部吃掉剩下的锥”的写法，与 \S3.4 的 $(K+\varepsilon H)$ 技巧同源，只是那里拆的是丰富性，这里拆的是数值等价。

（5）推论 6.41 把“模型有限”化归为“面有限”：任意被 $X$ 支配的 crepant 射影双有理 $Y\to Z$ 对应 $NE(X/Z)$ 的一个面 $NE(X/Y)$，而 $NE(X/Z)$ 只有有限多极射线，故只有有限多面。定理 6.42 [KM87] 再把它落到三维：对 $Z$ 上的 crepant 模型先作终端化与 Q-分解，化到终端情形后用推论 6.41。

（6）定理 6.44（标准翻模存在）对 $e(X)$ 作归纳：$D$ 已丰富时直接按定理 6.38 的 MMP 得翻模；否则先做一次 crepant 降下把 $D$ 搬到更小的模型上，$e(X)$ 严格下降；无处可降时由引理 6.39 判同构。定理 6.43 只证“射线型”标准翻模的终止，一般情形引 [Kol89] —— 这是全书少数\textbf{只引不证的主定理}，读时要标出来，不能当成已证。

（7）方法论收获：
\begin{itemize}
\item （1）“crepant 降下”= 把相对标准类为 0 的情形看成 $K+\varepsilon D$ 情形在 $\varepsilon\to0$ 且 $D$ 相对丰富时的极限；标准类为 0 不等于“没有正性可用”，正性藏在相对丰富的除子里。
\item （2）有限性定理（推论 6.41、定理 6.42）不只是“数模型”，它同时是终止性证明的\textbf{收网工具}：终止性常归结为“同构类有限，而每次翻模必改变某个不变量”。
\item （3）\S6.4 与 \S3.4 结构同源（定理 3.25(2) + 辅助丰富除子 + 有限生成），只是一个用来证有理性，一个用来数模型。
\end{itemize}

（8）新增疑问：
\begin{itemize}
\item （1）$e(X)$ 的严格下降在定理 6.44 证明的每一步怎么保证（书 p.254 只给归约框架，我还没完全还原归纳链）；
\item （2）[Kol89] 的一般终止性证明的核心不变量是什么（本书未展开）；
\item （3）术语陷阱：standard 指“$K_X$ 为 Q-Cartier 且奇点都是标准奇点”，不是“边界系数标准化”；我最初把“标准翻模”误当成“$\Delta$ 系数都被压到标准区间”，读到定理 6.44 的陈述才纠正。
\end{itemize}

\subsection*{九、我的三条方法论收获}

\begin{enumerate}
\item \textbf{"存在性"在双有理几何里通常是"某个环的有限生成性"}。第 6 章把翻转存在、Q-分解、终端化全部翻译成分次代数的有限生成性，再由有限生成性反推几何。这是 Mori--Kollár 学派最核心的操作：几何问题 $\leftrightarrow$ 代数（或"符号幂代数"）问题。
\item \textbf{局部化 + 奇点分类 = 存在的证明}。三维翻模存在性之所以可证，是因为三维终端奇点有显式局部模型；四维没有，所以方法断掉。\textbf{这告诉我：要看一个 MMP 定理能不能推到更高维，先问它的证明用了哪一维的哪条分类。}
\item \textbf{"不变量单调 + 边界分量逐个退出"是终止性证明的通用骨架}。三维用 $d(X,\Delta)$ 与后/前翻转曲线；四维以上 MMP 的终止性至今仍是开放问题，说明这只骨架在更高维还没有合适的"账本"。
\end{enumerate}

\subsection*{十、我的疑问}

\begin{itemize}
\item 引理 6.2 的 (1) 与 (2) 在\textbf{混合特征}下是否仍成立？注意 6.1 说定义与结果"仍有意义"，但 (2) 的"射影"在特征是 $p$ 时是否要换成"固有"？原书没有展开。
\item 定理 6.14 的证明只对 $K_X$ Cartier 情形完整；一般情形的"指数 1 覆盖上的循环群作用"细节书中没有给出，只提到分类（批注已记）。这是读第 6 章时最想补却补不上的一处。
\item 定理 6.17 的附加项（射线型标准 $(K_X+\Delta)$-翻转）与主项（$K_X$-翻转）之间，是否有一条统一的证明？书中是分开陈述的，我猜前者要靠"边界分量的数值独立"再加一层归纳，但没有验证。
\end{itemize}

\section*{研读札记 07 \textperiodcentered{} 第 7 章 半稳定极小模型（原书 p.257--298）}
\addcontentsline{toc}{section}{研读札记 07 \textperiodcentered{} 第 7 章 半稳定极小模型（原书 p.257--298）}
\markboth{研读札记 07 \textperiodcentered{} 第 7 章 半稳定极小模型（原书 p.257--298）}{}

第 7 章是全书的"收工章"：前面六章的工具（锥定理、无基点定理、终端化、Q-分解、翻模）在这里被组装成一台真正能跑的机器——半稳定三维 MMP。它也是全书唯一一章"先讲怎么退化、再讲翻转"的：把三维簇看成曲线上的一族曲面，翻转问题就变成"曲面族怎么坏掉、怎么修好"。

\subsection*{一、为什么要"半稳定"}

第 6 章已经证了三维终端翻模的存在性，但翻转的存在性（森 1988 / Sh92）证明很长很复杂。本章的做法是\textbf{换一个更受控的设定}：要求 $X\to C$ 的每个纤维都 lc。这样退化不会太坏，翻转可以在"看得见"的局部模型里逐纤维进行。

Kollár--Mori 在 [KM92] 里证明几乎所有翻转都是半稳定的——所以这不是回避，而是\textbf{把最难的部分自动满足掉}：你一旦知道"一般翻转都能扰动成半稳定"，剩下的就是把半稳定这一情形彻底搞清楚。

\subsection*{二、lc 射与 dlt 射：把"逐纤维"打包成一个条件（\S7.1）}

定义 7.1：$f:(X,B)\to C$ 是 \textbf{lc 射}，若对每个闭点 $c\in C$，$(X,B+f^{-1}(c))$ 是 lc；dlt 射同理。这个定义的妙处在于：它一次性把"每个纤维都 lc"和"$B$ 与纤维的相对位置"都锁定了。

引理 7.2 给了我五条"自动"：

\begin{enumerate}
\item 任意纤维约化——因为 $B+f^{-1}(c)$ 的系数 $\le1$。
\item $B$ 水平——即 $B$ 的不可约分支不含于纤维中。
\item $(X,B)$ 整体 lc。
\item 中心含于某纤维的例外除子 $E$：$a(E,X,B)\ge a(E,X,B+f^{-1}(c))+1\ge0$。\textbf{这条我读了两遍才明白它为什么"看起来像"标准对实际上不是}：它只说赋值非负，不说 $K_X+B$ 是 Q-Cartier。
\item \textbf{标准性/klt 性只取决于一般纤维} $X_{gen}$（把 $K_Z=h^*(K_X+B)+aE+\cdots$ 限制到 $X_{gen}$ 即得）。
\end{enumerate}

注 7.3 紧接着提醒：$(X,B)$ 一般不是标准的（反例 $X=C\times S$，$S$ lc 曲面）。但"表现得像标准对"——所以后面凡是需要"标准对"的定理，都值得试着改造成"lc 射"版本。\textbf{这一句是全章的方法论提示，我把它记下来了。}

引理 7.4 给了二维的局部形态：纤维要么在 $x$ 处是通常结点且 $x\notin\Supp B$，要么光滑且 $\mathrm{mult}_xB\le1$。引理 7.6 则说基变换保持 lc 射，证明链条是：平坦性 $\Rightarrow$ $S_2$ 保持 $\Rightarrow$ 只在余维 1 检查正规性；$f$ 在余维 2 之外光滑 $\Rightarrow$ $f'$ 的光滑集也能算；最后由 Hurwitz 公式 $K_{C'}+R'=g^*(K_C+R)$ 得 $K_{X'}+f'^{-1}(R')=h^*(K_X+f^{-1}(R))$。\textbf{这三步里我最想学的是"Hurwitz 公式把基曲线上的分歧搬到 $X$ 上"这一手。}

\subsection*{三、半稳定 MMP 的四个主定理（\S7.1）}

\begin{itemize}
\item \textbf{终止}（定理 7.7）：$(X,B)$ 三维 Q-分解 dlt、$g\circ f$ 为 lc 时，$Y$ 上任意一串射线型 $(K_X+B)$-翻转有限。
\item \textbf{翻转存在}（定理 7.8 $\to$ 推论 7.42）：同设定下翻转收缩有翻转（\S7.4 证）。
\item \textbf{MMP 存在}（定理 7.9）：在 $g\circ f$ 为 dlt 时，MMP 所有阶段存在且终止，终点满足 (1) 射仍 dlt、(2) $K_{X^m}+B^m$ 是 $f^m$-nef 或出现 Fano 纤维结构 $X^m\to Z^m\to Y$。证明就是 \S3.7 的运行加上"dlt 由推论 3.44 保持 + 翻转由推论 7.42 存在 + 终止由定理 7.7"。
\item \textbf{标准模型}（定理 7.10）：$K_X+B$ 是 $f$-大的时，在 $Y$ 上有标准模型，且射仍 lc。
\end{itemize}

有一点很值得注意：定理 7.10 在 dlt（非 klt）情形\textbf{只对三维成立}，因为要处理 $\lfloor B^m\rfloor$、而无基点定理不适用于这种情形，[K+92, 8.4] 用了曲面 $\lfloor B^m\rfloor$ 的详细分类结果。\textbf{这是"三维"这个限定在全书里最具体的一次出现。}

\subsection*{四、终止性的证明：把第 6 章的账本改一个过滤条件}

第 6 章终止性（定理 6.17）的证明几乎原封不动地搬过来：

\begin{itemize}
\item 定理 7.11 是局部版：设 $T\subset Y$ 是 $g$-平凡的 Cartier 除子、$S=f^*T$、$(X,S+B)$ lc，则 $Y\setminus T$ 上同构的 $(K_X+B)$-翻转序列有限。
\item 引理 7.12 是关键一步：\textbf{翻转收缩 $h$ 在 $g$ 的一般点附近必须是同构}。两个证明方法都很漂亮：情形 1 用刚性引理 1.6（$V_a$ 的整个纤维被压成一点，所以 $h$ 收缩的是一个除子，不可能是翻转收缩）；情形 2 用 $(-1)$-曲线的形变（把定理 1.10 阶段 2 的方法施于 $\ell'\to Z$，得形变空间维数 $\ge4$）。
\item 不变量换成定义 7.15 的 $d_S(X,\Delta)$：与 6.20 的 $d(X,\Delta)$ 唯一差别是\textbf{计数时只数中心含于 $S$ 的例外除子}。引理 7.14 保证 $\Sigma_S(0)=\emptyset$、klt 时 $\Sigma_S(1)$ 有限（用"分母是 $m$ 的约数"取 $\varepsilon<1/m$）。
\item 引理 7.16 就是 6.21 的对应版本；然后照 6.17 的归纳走完。
\end{itemize}

我在这里学到的：\textbf{终止性证明是"可移植"的，换设定只需要换不变量与局部引理，归纳骨架不变。} 这也解释了为什么三维 MMP 终止性在 1980 年代就能证，而四维以上至今没有——骨架有，账本没有。

\subsection*{五、半稳定约化（\S7.2）}

定理 7.17（[KKMSD73]）说：可以取有限基变换 + 射影消解，使 $f''$ 的纤维都是 snc 且约化。本章的贡献是定理 7.19：\textbf{在三维情形，不用环面嵌入的组合理论}，只用第 6 章的 Q-分解终端化就能得到半稳定奇点消解。

证明路线：先爆破成 $f^{-1}(c)+B$ snc；再取分歧覆盖 $C'\to C$ 使纤维约化（重数被分歧指数吃掉）；此时 $X'_n$ 只剩\textbf{标准商奇点}（命题 7.23 给局部模型），再用终端化 $X'_{tf}$ 把它们消掉，得到的 $X''$ 满足 snc 结论。注 7.20 说二维版本就是 [DM69]，那里"非奇异射影终端化"就是极小奇点消解。

\textbf{这里我看到"第 6 章的工具第一次被当作黑箱使用"}：终端化定理 6.23 与 Q-分解定理 6.25 在 \S7.2 里只是两个"取一下就有"的操作，而它们的证明曾经占了第 6 章一半篇幅。这是全书结构上最漂亮的一处：先造工具，再用工具。

\subsection*{六、特殊半稳定翻转（\S7.3）}

定理 7.32 是三假设下的存在性：$(X,S+B_X)$ dlt、$(X\setminus(S\cup\lfloor B_X\rfloor),0)$ 终端、$S$ 有分支 $S_1$ 与 $C=\mathrm{Ex}(f)$ 负交。结论：$f$ 有翻转。

证明的核心是命题 7.34：\textbf{造出一个 lc 补除子}，把翻转问题变成三维标准翻模的存在问题。这样就把 \S7.3 归约到第 6 章的定理 6.14 上了——\textbf{"用翻模造翻转"}，这是本章最反直觉的一步：翻转（$K$ 为负）居然要靠翻模（$K$ 平凡）来造。

命题 7.38 是纯曲面（二维）的整理工具：$p:S\to T$ 奇点消解、$(S,D+\Delta)$ plt、$-(K_S+D+\Delta)$ 是 $p$-nef，则存在 $D^\sharp_S$ 使 $(S,D+\lfloor2\Delta\rfloor+D^\sharp_S)$ 是 lc。它把"负数"整理成"lc 的边界"，是 \S7.3 的技术底座。

\subsection*{七、一般半稳定翻转（\S7.4）}

主定理 7.40：三维 klt 对、$s\circ f$ 为 lc，则翻转收缩有翻转。推论 7.42 把它升到 dlt：把 $B$ 换成 $(1-\varepsilon)B$ 不改变翻转（定义 6.5），于是归约到 klt。

证明（7.43）分两阶段：\textbf{“减去 $H$”（theorem 7.44）} 与 \textbf{定理 7.51}，中间夹一个"容易阶段 7.54"。思想是 [Sho92]、这里的论证基于 Corti：先作一个巧妙的辅助设定，使该设定下的 MMP \textbf{只需要特殊半稳定翻转与除子收缩}就能跑完——于是 \S7.3 就够了，不需要一般翻转。

这一步我理解为"把未知的翻转换成已知的翻转 + 除子收缩"：MMP 里最危险的步骤（翻转）被限制在特殊情形，其余都由除子收缩完成。\textbf{这正是"用 MMP 证明 MMP"的自举结构，也是全书方法论的高峰。}

\subsection*{八、应用（\S7.5）}

三个应用共用同一台机器：

\begin{enumerate}
\item \textbf{一般型曲面模空间的紧化}：稳定曲线（定义 7.57：$g$ 维固有、奇点仅通常二重点、$\omega_D$ 丰富、$\chi=1-g$）与稳定曲线族，把 $\overline{M}_g$ 的构造推广到曲面。
\item \textbf{曲面奇点的半普遍形变空间}。
\item \textbf{曲面族的同时奇点消解}（定理 7.51 给出条件的充要性）。
\end{enumerate}

\S7.5 的详细推导我没有逐行读，但结构很清楚：\textbf{半稳定 MMP 是一台"把曲面族的坏退化变成好退化"的机器}，凡是需要"同时"修好一族曲面（同时消解、模空间紧化）的问题，都可以把族拉到半稳定 MMP 的终点，靠终点纤维的 snc 性来做事。

\subsection*{九、\S7.6 文献概观：一张 1998 年的地图}

（1）定位。\S7.6 不证明任何东西，是“文献表 + 一句话状态”。它的价值不在结论，而在给出每个方向的\textbf{边界}：哪些已解决、哪些只对特殊维数解决、哪些完全未解。我把它当索引用，逐条对照正文方法。

（2）七条线的骨架：
\begin{itemize}
\item \textbf{翻转（7.71）}：3 维终端情形 [Mor88]，3 维 lc 情形 [Sho92]／[K+92]；\textbf{4 维以上未知}，书中直言“这大概是极小模型理论中最重要的未解决问题”，已知只有 toric [Rei83a] 与特殊群作用 [BK94]。这一句直接标定本书的历史位置：三维做完了，四维打开。
\item \textbf{翻转与 GIT（7.72）}：极化不同则商空间不同，但彼此双有理等价，其间映射往往由 flip 或逆 flip 复合而成（[Tha96, DH98, MW97]）。这解释了“flip 不是人为构造，而是 GIT 变的必然产物”。
\item \textbf{小平维数（7.73）}：$\kappa(X)$ 的定义与饭高茂的起点；[Uen75]、[Mor87]。
\item \textbf{丰度与标准环（7.74）}：标准环有限生成性；3 维 $\kappa\le2$ 由 [Fuj86]，3 维丰度由川又--宫冈系列与 [KMM94a]；另有“常数 $c(n)$”问题：$\kappa=0$ 由 [Kaw86]，$\kappa=1$ 是森（未发表），$\kappa=2$ 由 [Kol94a]，$\kappa=3$ 未解决。这条最实用：它把“丰度猜想 = 标准环有限生成”的等价链边界说清楚了。
\item \textbf{藤田猜想（7.75）}：$m\ge\dim X+1$ 时 $|K_X+mL|$ 自由；$\dim\le4$ 成立 [Kaw97]，一般维数下 $m>\frac12\dim X(\dim X+1)$ 时自由 [AS95]；副产品：不可分解主极化 Abel 簇 $(A,\Theta)$ 是标准的（$\Theta$ 正规、只有有理奇点，[EL97]）。
\item \textbf{Fano 纤维空间的双有理几何（7.76）}：与标准模型唯一相对照，Fano 终点\textbf{不唯一}，判别双有理等价极难；两条路线——选特殊边界（[Ale91, Cor96, CF93]）与 Sarkisov 纲领式分解（[Cor94]）。这正好解释第 3 章标准模型唯一性定理为何重要：唯一性是特例，不唯一才是常态。
\item \textbf{三维 Fano／Calabi--Yau／椭圆簇（7.77--7.79）}：非奇异 3 维 Fano 有 104 个形变型 [Isk80, MM81]；$\rho=1$ 的 Q-分解终端 3 维 Fano 至多有限多形变型 [Kaw92a]，标准情形 [KMM94b] 且无需 $\rho=1$；3 维 CY 的形变型可能无限（[Rei87a] 猜想、[NS95] 光滑化）；3 维椭圆簇 [Nak91, DG94]。
\item \textbf{复簇与正标数（7.80--7.81）}：固有／紧复簇上曲线锥不再是正确对象（有效曲线可数值平凡），Kähler 情形反而可行 [CP97]；正标数／混合标数的 MMP 在 1998 年仍是信念而非定理。
\end{itemize}

（3）方法论收获：
\begin{itemize}
\item （1）\textbf{“概说”节是作者最诚实的自留地}：哪条完整证明、哪条引用、哪条是猜想，这里分得最清。我以后拿不准“这是定理还是期待”时，先回 \S7.6 对表。
\item （2）文献表本身就是研究地图：7.71 的四维翻转、7.74 的 $c(n)$、7.76 的 Fano 双有理分类，是三个至今仍活跃的方向。
\item （3）用词纪律：\textbf{丰度（abundance）$\ne$ 丰富（ample）}；\textbf{crepant 是相对于对 $(X,\Delta)$ 而言}；\textbf{标准（canonical）指奇点类别，不是系数}。
\item （4）\S7.6 的 7.71--7.82 是“编号段”不是定理，尾符统一 $\square$；若误当定理编号，会以为第 7 章有二十多个定理。
\end{itemize}

（4）新增疑问：
\begin{itemize}
\item （1）7.74 说“假定 MMP 则丰度 $\Rightarrow$ 有限生成，但逆大概不成立”，逆不成立的反例是什么？
\item （2）7.77“104 个形变型”的计数口径（是否按形变族、是否含 $K$-等价重复）？
\item （3）7.80 说曲线锥“本来就不正确”，那么复解析情形替代它的是什么对象（Moishezon／Kähler 锥？），本书未展开。
\end{itemize}

\subsection*{十、我的四条方法论收获}

\begin{enumerate}
\item \textbf{"特殊设定 + 归约定理"比"一般证明"更有力}。半稳定翻转比一般翻转好证得多，而 [KM92] 保证一般情形几乎都能半稳定化。这是数学里常见的策略：把难的陈述化成"一个特殊情形 + 一个归约"。
\item \textbf{把逐纤维条件打包成一个射的条件}（定义 7.1）能让"保持性"变得可证：推论 3.44（dlt 保持）、引理 7.6（基变换保持 lc）都是这种打包的直接红利。
\item \textbf{终止性证明的骨架可以跨设定移植}：不变量（$d$ 或 $d_S$）+ 局部引理（6.21/7.16）+ 对边界分量个数归纳。三维能跑，四维不能跑，差别在"账本"是否有限。
\item \textbf{工具的价值在第二章才显现}。第 6 章的终端化/Q-分解在 \S7.2 里变成一句话的引用；这是"先造工具再使用"的教科书式示范，也解释了为什么作者要把第 6 章写得那么细。
\end{enumerate}

\subsection*{十一、我的疑问}

\begin{itemize}
\item 定理 7.10 的 dlt 情形只对三维成立（依赖曲面 $\lfloor B^m\rfloor$ 的结果）。那么四维的半稳定标准模型定理目前是什么状态？书中没有说。
\item 定理 7.32 的假设 (2)（$(X\setminus(S\cup\lfloor B_X\rfloor),0)$ 终端）是否是自动的？若 $S$ 是 $f$-平凡的有效 Cartier 除子且 $(X,S+B_X)$ dlt，直觉上"离开 $S$ 之后应当是终端的"，但书中没有说明，我怀疑这正是"特殊"二字的来源。
\item 注 7.5 提到 slc 对（半对数标准）可以处理非正规奇点，但没有展开。slc 在 \S7.6 的"最近发展"里是否被用到？要等 \S7.6 读完再回来核。
\end{itemize}
